\chapter{Financial Statements \& Accounting}\label{ch:financial-statements}

% Scope: the three-statement model & linkages, accrual vs cash (deferred revenue),
%   cost accounting (contribution margin vs gross margin, ABC), the statement-to-ROIC bridge.
%   NEW chapter (closes the highest-impact gap: ch.19--22 use NOPAT/FCF/EBITDA undefined).
% Prerequisite for \cref{ch:corporate-finance-core}. PINS the cost structure that yields the
%   book's 42% contribution margin / $21/mo --- reconcile, do not contradict, prior uses.
% Draft per house style: flowing prose + example/admonition set-offs; running-SaaS worked examples.

Every valuation number in this book --- the \money{600000} NOPAT, the \money{3000000}
invested capital, the \qty{20}{\percent} ROIC, the \money{3150} CLV --- is derived from
financial statements. Operators who cannot trace those numbers back to income statements,
balance sheets, and cash-flow statements are working from faith, not analysis. This chapter
builds the accounting substrate that the rest of the book assumes: how the three statements
link, why accrual and cash diverge, how cost structure determines the contribution margin
that governs pricing and acquisition economics, and how to convert statutory accounting into
the NOPAT and free-cash-flow inputs that feed discounted-cash-flow valuation.

\section{The Three-Statement Model}\label{sec:three-statements}
% complexity: 4

How does a sale today flow through a company's financial position over the next twelve
months? Most business-school graduates can name the three financial statements; fewer can
trace a single transaction across all three simultaneously, and fewer still understand why
that linkage matters for valuation. The income statement records revenues and expenses over
a period. The balance sheet records what the firm owns and owes at a point in time. The cash
flow statement reconciles the two by tracking the actual movement of cash. None of the three
is interpretable in isolation; all three are required for a complete picture.

The architecture of the linkage is precise. Net income --- the bottom line of the income
statement --- flows directly into retained earnings, the cumulative equity account on the
balance sheet; every dollar of profit not paid out as a dividend accumulates there. The cash
flow statement then starts with net income and adjusts it for two systematic distortions:
non-cash charges (principally depreciation and amortisation, which reduce income but consume
no cash) and changes in working capital (receivables, payables, and inventory that shift the
timing of cash relative to revenue recognition). Capital expenditure flows from the investing
section of the cash flow statement into property, plant, and equipment on the balance sheet,
increasing the asset base that generates future income. Debt raised or repaid appears in the
financing section and directly changes the liability side of the balance sheet.

These four linkages define the cash identity across any period:
\begin{equation}\label{eq:three-statement-link}
  \text{Cash}_t \;=\; \text{Cash}_{t-1} \;+\; \mathrm{CFO}_t \;+\; \mathrm{CFI}_t \;+\; \mathrm{CFF}_t,
\end{equation}
where CFO is cash from operations (net income adjusted for D\&A and working capital changes),
CFI is cash from investing (primarily capital expenditure, negative when spending), and CFF
is cash from financing (debt issuance positive, repayment negative). The constraint is
accounting arithmetic: the three activity categories must sum to the change in cash,
\(\text{Cash}_t - \text{Cash}_{t-1}\). A company can run negative CFO indefinitely only if
CFF --- external capital from investors or lenders --- is large enough to compensate. This is
precisely the growth-phase condition that matters for early-stage SaaS firms.

\begin{figure}[htbp]
  \centering
  \includegraphics[width=0.9\linewidth]{three-statement}
  \caption{The three-statement linkage. Net income flows to retained earnings (equity);
    D\&A is added back in CFO; CapEx reduces CFI and increases PP\&E (net of depreciation);
    debt proceeds flow through CFF to the liability side of the balance sheet.
    \Cref{eq:three-statement-link} captures the resulting cash identity.}
  \label{fig:three-statement}
\end{figure}

\begin{example}[Year-1 SaaS: growth-phase cash burn]
The running firm ends Year~0 with \num{6667} paying customers (ARR \money{4000000} at
\money{600}/customer/year). In Year~1 it adds \num{2400} new customers (\num{200}/month at
\money{600} CAC), ending the year at \num{9067} customers and ARR of roughly \money{5440000}.
Tracing through the three statements:

\medskip
\textit{Income Statement.} Annual revenue \money{4000000} (Year-1 cohort; new customers
contribute a partial year). Variable COGS at \qty{30}{\percent} of revenue (\money{15}/month
per customer) yields gross profit \money{2800000} (\qty{70}{\percent} gross margin). Fixed
operating costs --- engineering, G\&A, product --- total roughly \money{1940000}. EBIT is
therefore \money{860000} and, taxed at \qty{30}{\percent}, net income is approximately
\money{602000}.

\medskip
\textit{Balance Sheet.} Net income \money{602000} increases retained earnings. The firm
spends \money{300000} on servers and software infrastructure (CapEx), increasing PP\&E by
that amount net of \money{60000} depreciation (\money{240000} net addition). It also invests
\money{1440000} in customer acquisition (\num{2400} customers at \money{600} CAC), funded
from operating cash.

\medskip
\textit{Cash-Flow Statement.} CFO begins with net income \money{602000}; add back D\&A
\money{60000}; subtract the increase in accounts receivable from the growing customer base
(\money{40000}), yielding CFO of approximately \money{622000}. CFI is \money{-300000}
(infrastructure CapEx). Customer acquisition spend of \money{1440000} is an operating cash
outflow that sits within CFO (sales and marketing expense). Applying \cref{eq:three-statement-link},
the net cash movement is \money{622000} + \money{-300000} = \money{322000} before
acquisition spend; after the \money{1440000} CAC outflow, the firm is net-cash-negative by
approximately \money{1118000} on the year.

\medskip
The key implication: the income statement reports \money{602000} of net income while the
firm actually consumed over \money{1000000} of cash. The business is accrual-profitable in
its steady-state NOPAT sense --- \cref{sec:statement-analysis} will derive that
\money{600000} precisely --- but cash-negative in Year~1 because investment in new customers
(a cash outflow today) precedes the revenue those customers generate (a cash inflow spread
across future years). Growth-phase cash burn is not a sign of a broken business; it is the
arithmetically inevitable consequence of front-loaded acquisition spend. The three-statement
model makes this visible; P\&L alone obscures it.
\end{example}

\begin{admonition}{Note}
The D\&A add-back in CFO is often misread as ``money the company gets back.'' It is not. D\&A
is a non-cash charge --- the cash went out when the asset was purchased (recorded as CFI at
that time). The add-back simply reverses an accounting deduction that reduced net income
without consuming new cash in the period. Confusing the add-back with free cash is one of
the more common errors in first-pass financial analysis.
\end{admonition}

The three-statement architecture is the skeleton on which all the chapters that follow are
built. \textcite{libby2022financial} provide the canonical treatment of transaction analysis
and the building-block approach that traces individual entries across all three statements.
\textcite{weygandt2019financial} develop the same framework with detailed attention to the
indirect method of the cash-flow statement used by most public companies. The investment
implications of the linkage --- and the derivation of free cash flow from these statements
--- are developed in \textcite{brealey2020}.

\section{Accrual Accounting vs.\ Cash}\label{sec:accrual}
% complexity: 3

When does revenue exist? The answer depends entirely on which accounting system you are
using, and the choice has material consequences for any SaaS company that collects
subscriptions in advance. Under \emph{cash accounting}, revenue is recognised the moment
cash is received. Under \emph{accrual accounting} --- the system required by GAAP and IFRS
for any company of meaningful size --- revenue is recognised when the service obligation is
fulfilled, regardless of when cash changes hands. The gap between the two shows up on the
balance sheet as deferred revenue: a liability representing cash received for services not
yet delivered.

For a subscription business with annual prepayment, the identity at any point in the year is:
\begin{equation}\label{eq:accrual-revenue}
  \text{Cash received} \;=\; \text{Revenue recognised} \;+\; \Delta\text{Deferred Revenue},
\end{equation}
where \(\Delta\text{Deferred Revenue}\) is the increase in the deferred revenue liability
over the period. On January~1, annual cash collected equals recognised revenue (zero days
elapsed) plus the full amount deferred; at December~31 the deferred balance is zero and all
revenue has been recognised. The equation holds at every moment in between. \Cref{eq:accrual-revenue}
is an accounting identity, not a model --- it cannot be violated.

\begin{example}[Annual prepayment: cash vs.\ P\&L]
On 1~January, \num{100} customers each prepay \money{600} for a year of software access.
Cash received is \money{60000}. Under accrual accounting (\cref{eq:accrual-revenue}),
January's P\&L may recognise only one month of service delivered:
\[
  \text{Revenue recognised in January} \;=\; \frac{\money{60000}}{12} \;=\; \money{5000}.
\]
The remaining \money{55000} sits on the balance sheet as deferred revenue --- a liability,
because the firm still owes eleven months of service. Applying \cref{eq:accrual-revenue}:
\[
  \money{60000} \;=\; \money{5000} \;+\; \money{55000}. \checkmark
\]
By 31~December, the deferred revenue balance has been drawn to zero and all \money{60000} has
been recognised as revenue. Cash never changed again; only the accounting classification
shifted. The annual P\&L and the cash-flow statement agree at year-end, but monthly P\&Ls
look radically different from monthly cash flows throughout the year --- a distinction that
matters for covenants, board reporting, and any investor reading a mid-year snapshot.

The mechanics here are the same as those governing MRR and ARR in \cref{sec:saas-mrr}:
annual prepayment is a form of working-capital financing by customers, and the deferred
revenue balance is the firm's floating obligation to deliver future service.
\end{example}

\begin{admonition}{Warning}
Annual prepay is a source of \emph{favourable} working capital for a SaaS firm: customers
finance operations in advance. That advantage disappears immediately if the company shifts to
monthly billing or if churn accelerates before the obligation is fulfilled. Investors reading
a SaaS balance sheet should treat a large deferred revenue balance as a quality signal only
when retention is high; a large deferred balance alongside rising churn is a future revenue
shortfall hiding as a present liability.
\end{admonition}

The revenue-recognition rules governing subscription businesses were formalised under ASC~606
(US~GAAP) and IFRS~15 (international), both of which shifted from a rules-based framework to
a principles-based performance-obligation model. For multi-element SaaS contracts ---
implementation services bundled with recurring licences --- the allocation of transaction
price across obligations requires judgment that can materially affect the timing of revenue
recognition. \textcite{weygandt2019financial} and \textcite{libby2022financial} both cover
the ASC~606 five-step model; the working-capital implications for subscription operators are
developed further in the operating model of \cref{ch:operating-model}.

\section{Cost Accounting: Contribution Margin \& Cost Allocation}\label{sec:cost-accounting}
% complexity: 3

How much does one more customer actually cost, and how much does one more customer actually
contribute? Gross margin and contribution margin both claim to answer this question, and they
will give different numbers for the same business. Understanding which number to use --- and
when --- is one of the most practically important distinctions in managerial accounting.

Gross margin is a statutory, GAAP-regulated metric: revenue minus the cost of goods sold
(COGS), where COGS includes both variable production costs and any fixed overhead allocated
to each unit. For a SaaS firm, COGS is primarily hosting, infrastructure, and direct
technical support costs --- the variable fraction of serving each customer. Contribution
margin is a non-GAAP managerial metric: revenue minus \emph{all} variable costs, whether
those costs live in COGS, sales, or operations. The distinction matters because SaaS firms
carry variable costs outside COGS --- payment processing fees, customer success labour tied
to customer count, variable support costs --- that are real economic costs of each incremental
customer but invisible in the gross margin line.

The contribution margin equation is:
\begin{equation}\label{eq:contribution-margin}
  \mathrm{CM} \;=\; \text{Revenue} \;-\; \text{Total variable costs},
  \qquad
  \mathrm{CM\%} \;=\; \frac{\mathrm{CM}}{\text{Revenue}}.
\end{equation}
This is the floor that must be cleared before fixed costs contribute to profit. Every unit
sold at a price above its variable cost contributes something to covering fixed overhead and
eventually generating NOPAT; pricing below variable cost destroys value on every unit regardless
of volume.

For a more precise cost allocation --- particularly when multiple product lines or customer
segments share infrastructure and support resources --- Activity-Based Costing assigns costs
to products in proportion to their actual consumption of each overhead activity rather than
spreading overhead evenly by headcount or revenue. An illustrative ABC unit cost for a given
customer segment is:
\begin{equation}\label{eq:abc-cost}
  \text{Unit cost}_{\text{ABC}} \;=\; \sum_{a} \frac{\text{Cost pool}_a}{\text{Activity volume}_a} \times \text{Activity consumed by unit},
\end{equation}
where each cost pool \(a\) represents a distinct overhead activity (server provisioning,
support ticket resolution, payment processing). ABC prevents high-volume simple products from
subsidising low-volume complex ones, which conventional gross-margin analysis systematically
obscures \parencite{datar2025horngrens}.

\begin{example}[Reconciling 70\% gross margin and 42\% contribution margin]
The firm charges \money{50}/month per customer. Variable COGS --- hosting and infrastructure
costs that scale directly with customer count --- are \money{15}/month per customer:
\[
  \text{Gross margin} \;=\; \frac{\money{50} - \money{15}}{\money{50}} \;=\; \frac{\money{35}}{\money{50}} \;=\; \qty{70}{\percent}.
\]
But the firm also carries additional variable costs outside COGS: customer support staff
(variable with ticket volume, which scales with customers), payment-processing fees (a
percentage of each transaction), and customer-success coverage (headcount tied to customer
count). These total approximately \money{14}/month per customer. Total variable cost per
customer per month is therefore \money{15} + \money{14} = \money{29}. Applying
\cref{eq:contribution-margin}:
\[
  \mathrm{CM} \;=\; \money{50} \;-\; \money{29} \;=\; \money{21}/\text{month},
  \qquad
  \mathrm{CM\%} \;=\; \frac{\money{21}}{\money{50}} \;=\; \qty{42}{\percent}.
\]
Annually, each customer contributes \money{21} \(\times\) 12 = \money{252}/year to covering
fixed costs and generating NOPAT.

The implication for pricing: the firm must charge above \money{29}/month on every plan to
avoid destroying value. The \money{50} price clears that floor by \money{21} --- a real but
not extravagant margin. A competitor who undercuts to \money{35}/month believing gross margin
is the binding constraint (\money{35} \(>\) \money{15} COGS, so apparently safe) is actually
pricing below total variable cost and losing \money{6}/month per customer at scale.
\end{example}

\begin{admonition}{Warning}
Contribution margin is not gross margin, and conflating them produces systematically wrong
pricing and acquisition decisions. Gross margin (\qty{70}{\percent}) counts only variable
COGS against revenue; contribution margin (\qty{42}{\percent}) counts \emph{all} variable
costs. The \qty{28}{\percent} gap --- \money{14}/month per customer in this firm --- represents
real economic costs that the gross margin line simply does not capture. A CAC payback
calculation run on gross margin overstates the rate at which acquisition cost is recovered;
the correct payback uses contribution margin. \Cref{eq:cac-payback} in the customer-economics
chapters implements this correctly.
\end{admonition}

The ``different costs for different purposes'' principle \parencite{datar2025horngrens}
applies directly: gross margin is the right metric for statutory reporting and for comparing
infrastructure efficiency across SaaS peers; contribution margin is the right metric for
pricing, acquisition decisions, and break-even analysis. \textcite{weygandt2019financial}
cover the GAAP treatment of COGS; the managerial decision-making applications of contribution
margin analysis are developed in detail in \textcite{datar2025horngrens}.

\section{Statement Analysis: From Accounting to ROIC}\label{sec:statement-analysis}
% complexity: 4

Financial statements serve statutory reporting purposes; valuation requires economic measures.
The transition from one to the other is not cosmetic. GAAP net income is distorted by
financing choices (interest deductions), tax effects, and non-cash accounting entries.
Enterprise value rests on the operating cash flows that belong to \emph{all} capital
providers, independent of how the firm is financed. Converting accounting to economics
requires two steps: isolating unlevered operating profit (NOPAT) and then adjusting for the
cash that investing activities consume or release (FCF).

NOPAT strips out the effect of financial leverage to reveal what the business earns on its
operating assets. FCF then adjusts NOPAT for the non-cash depreciation charge (add back,
since the cash left when the asset was bought, not when it depreciates) and for the cash
consumed by new capital investment (CapEx) and changes in working capital (\(\Delta\)WC).
The bridge is:
\begin{equation}\label{eq:nopat-bridge}
  \mathrm{NOPAT} \;=\; \mathrm{EBIT} \times (1 - t),
  \qquad
  \mathrm{FCF} \;=\; \mathrm{NOPAT} \;+\; \mathrm{D\&A} \;-\; \Delta\mathrm{WC} \;-\; \mathrm{CapEx},
\end{equation}
where \(t\) is the effective tax rate, \(\Delta\mathrm{WC}\) is the increase in net working
capital (positive when the business is consuming cash to fund growth), and CapEx is gross
capital expenditure on long-lived assets. ROIC is NOPAT divided by invested capital (IC),
the total capital --- equity plus interest-bearing debt --- deployed to generate that profit:
\(\mathrm{ROIC} = \mathrm{NOPAT} / \mathrm{IC}\). When ROIC exceeds the cost of that
capital (\cref{eq:economic-profit}), the firm creates economic value; when it falls below,
it destroys it, regardless of what net income reports.

\begin{example}[Deriving NOPAT, FCF, and ROIC from the income statement]
From the three-statement example above, the firm's income statement shows EBIT of
\money{860000} at the \num{6667}-customer steady state (ARR \money{4000000}, variable COGS
\money{1200000} at \qty{30}{\percent} of revenue, additional variable opex \money{1120000},
fixed opex \money{820000} including \money{60000} D\&A). Applying \cref{eq:nopat-bridge}
with a \qty{30}{\percent} tax rate:
\[
  \mathrm{NOPAT} \;=\; \money{860000} \times (1 - 0.30) \;=\; \money{860000} \times 0.70 \;=\; \money{602000} \;\approx\; \money{600000}.
\]
This is the \money{600000} NOPAT referenced throughout the book.

To derive FCF, add back D\&A (\money{60000}) and assume steady-state CapEx of \money{60000}
(replacement only, so CapEx equals D\&A) and working-capital change \(\Delta\mathrm{WC} = 0\)
at steady state (new-customer cash is collected monthly, no meaningful receivables build):
\[
  \mathrm{FCF} \;=\; \money{600000} \;+\; \money{60000} \;-\; \money{0} \;-\; \money{60000} \;=\; \money{600000}.
\]
Invested capital is the equity deployed to fund the business at steady state --- servers,
software development, and the working capital needed to operate: IC = \money{3000000}.
\[
  \mathrm{ROIC} \;=\; \frac{\money{600000}}{\money{3000000}} \;=\; \qty{20}{\percent}.
\]
WACC is \qty{11}{\percent} (\cref{eq:wacc}), so the ROIC--WACC spread is \qty{9}{\percent}.
From \cref{eq:economic-profit}, economic profit is \(\money{3000000} \times 0.09 =
\money{270000}\)/year --- meaning the business generates \money{270000}/year of value above
and beyond the return required by its capital providers.

These are the numbers \cref{eq:roic} and \cref{eq:dcf} use throughout. The \money{600000}
FCF, discounted at WACC with growth \(g = \qty{3}{\percent}\) over five explicit years plus
a terminal value, yields the enterprise value developed in the DCF chapter (\cref{eq:dcf},
\cref{eq:npv}). Every link in that chain starts here, in the income statement.
\end{example}

\begin{admonition}{Warning}
EBITDA is not free cash flow. EBITDA adds D\&A back to EBIT but makes no deduction for CapEx
or changes in working capital. A capital-intensive business can report strong EBITDA while
consuming nearly all of it in maintenance CapEx; its FCF is a fraction of EBITDA. For SaaS
firms with modest physical CapEx, the difference is smaller but not zero: working capital
build and intangible development costs can still create a material gap. Lenders and acquirers
who rely on EBITDA multiples without FCF conversion are systematically overpaying for
capital-intensive businesses and underpaying for capital-light ones. The correct valuation
input is always FCF derived from \cref{eq:nopat-bridge}, not EBITDA.
\end{admonition}

The NOPAT--FCF--ROIC bridge developed here feeds directly into the DCF model of
\cref{eq:dcf} and the economic-profit test of \cref{eq:economic-profit}. Investors use
\cref{eq:roic} to assess whether the firm earns its cost of capital; operators use the same
calculation to evaluate capital-allocation decisions. \textcite{koller2025valuation} provide
the definitive treatment of this conversion and its applications to enterprise valuation.
\textcite{brealey2020} connect the accounting bridge to the NPV rule of \cref{eq:npv} and
the weighted-average cost of capital of \cref{eq:wacc}.
